\documentclass[10pt,a4paper]{amsart}
\usepackage{amsmath,amssymb,amsthm,graphicx}
% Retained preamble for the Stacks Project chapter snapshot in this directory.
% The chapter's own preamble.tex is a copy of the shared Stacks Project preamble; the block
% below is sliced MECHANICALLY, line for line, from it. Nothing was added, removed or
% redefined; the wrapper class, the xr-hyper/hyperref wiring and the externaldocument list
% are omitted because the assembler supplies the document class.
%
% lines taken from preamble.tex: 3, 5, 9, 11, 12, 156, 157, 158, ... (51 lines total)
%
% package coverage restored from the chapter's own preamble (2026-10-13): xy, xy-2cell

\usepackage[all]{xy}\usepackage{verbatim}
\xyoption{2cell}
\UseAllTwocells

\newenvironment{reference}{\comment}{\endcomment}
\newenvironment{slogan}{\comment}{\endcomment}
\newenvironment{history}{\comment}{\endcomment}
\theoremstyle{plain}
\newtheorem{theorem}[subsection]{Theorem}
\newtheorem{proposition}[subsection]{Proposition}
\newtheorem{lemma}[subsection]{Lemma}
\theoremstyle{definition}
\newtheorem{definition}[subsection]{Definition}
\newtheorem{example}[subsection]{Example}
\newtheorem{exercise}[subsection]{Exercise}
\newtheorem{situation}[subsection]{Situation}
\theoremstyle{remark}
\newtheorem{remark}[subsection]{Remark}
\newtheorem{remarks}[subsection]{Remarks}
\numberwithin{equation}{subsection}
\def\lim{\mathop{\mathrm{lim}}\nolimits}
\def\colim{\mathop{\mathrm{colim}}\nolimits}
\def\Spec{\mathop{\mathrm{Spec}}}
\def\Hom{\mathop{\mathrm{Hom}}\nolimits}
\def\Ext{\mathop{\mathrm{Ext}}\nolimits}
\def\SheafHom{\mathop{\mathcal{H}\!\mathit{om}}\nolimits}
\def\SheafExt{\mathop{\mathcal{E}\!\mathit{xt}}\nolimits}
\def\Sch{\mathit{Sch}}
\def\Mor{\mathop{\mathrm{Mor}}\nolimits}
\def\Ob{\mathop{\mathrm{Ob}}\nolimits}
\def\Sh{\mathop{\mathit{Sh}}\nolimits}
\def\NL{\mathop{N\!L}\nolimits}
\def\CH{\mathop{\mathrm{CH}}\nolimits}
\def\proetale{{pro\text{-}\acute{e}tale}}
\def\etale{{\acute{e}tale}}
\def\QCoh{\mathit{QCoh}}
\def\Ker{\mathop{\mathrm{Ker}}}
\def\Im{\mathop{\mathrm{Im}}}
\def\Coker{\mathop{\mathrm{Coker}}}
\def\Coim{\mathop{\mathrm{Coim}}}
\DeclareMathSymbol{\boxtimes}{\mathbin}{AMSa}{"02}
\def\QCohstack{\mathcal{QC}\!\mathit{oh}}
\def\Cohstack{\mathcal{C}\!\mathit{oh}}
\def\Spacesstack{\mathcal{S}\!\mathit{paces}}
\def\Quotfunctor{\mathrm{Quot}}
\def\Hilbfunctor{\mathrm{Hilb}}
\def\Curvesstack{\mathcal{C}\!\mathit{urves}}
\def\Polarizedstack{\mathcal{P}\!\mathit{olarized}}
\def\Complexesstack{\mathcal{C}\!\mathit{omplexes}}
\def\Pic{\mathop{\mathrm{Pic}}\nolimits}
\def\Picardstack{\mathcal{P}\!\mathit{ic}}
\def\Picardfunctor{\mathrm{Pic}}
\def\Deformationcategory{\mathcal{D}\!\mathit{ef}}
\usepackage{cleveref}
\crefname{originallemma}{original lemma}{original lemmas}
\Crefname{originallemma}{Original lemma}{Original lemmas}
\makeatletter
\newlabel{definition-approximation}{{08HJ}{source}{}{sourcepointer-1}{}}
\newlabel{definition-approximation@cref}{{[originallemma][0][]08HJ}{[1][0][]source}{}{}{}}
\makeatother
\crefname{originallemma}{original lemma}{original lemmas}
\Crefname{originallemma}{Original lemma}{Original lemmas}
\makeatletter
\newlabel{definition-approximation-holds}{{08HI}{source}{}{sourcepointer-2}{}}
\newlabel{definition-approximation-holds@cref}{{[originallemma][0][]08HI}{[1][0][]source}{}{}{}}
\makeatother
\crefname{originallemma}{original lemma}{original lemmas}
\Crefname{originallemma}{Original lemma}{Original lemmas}
\makeatletter
\newlabel{derived-lemma-cohomology-of-hocolim}{{0CRK}{source}{}{sourcepointer-3}{}}
\newlabel{derived-lemma-cohomology-of-hocolim@cref}{{[originallemma][0][]0CRK}{[1][0][]source}{}{}{}}
\makeatother
\crefname{originallemma}{original lemma}{original lemmas}
\Crefname{originallemma}{Original lemma}{Original lemmas}
\makeatletter
\newlabel{derived-lemma-representable-homological}{{0149}{source}{}{sourcepointer-4}{}}
\newlabel{derived-lemma-representable-homological@cref}{{[originallemma][0][]0149}{[1][0][]source}{}{}{}}
\makeatother
\crefname{originallemma}{original lemma}{original lemmas}
\Crefname{originallemma}{Original lemma}{Original lemmas}
\makeatletter
\newlabel{lemma-ext-from-perfect-into-bounded-QCoh}{{09MB}{source}{}{sourcepointer-5}{}}
\newlabel{lemma-ext-from-perfect-into-bounded-QCoh@cref}{{[originallemma][0][]09MB}{[1][0][]source}{}{}{}}
\makeatother
\crefname{originallemma}{original lemma}{original lemmas}
\Crefname{originallemma}{Original lemma}{Original lemmas}
\makeatletter
\newlabel{sites-cohomology-definition-pseudo-coherent}{{08FT}{source}{}{sourcepointer-6}{}}
\newlabel{sites-cohomology-definition-pseudo-coherent@cref}{{[originallemma][0][]08FT}{[1][0][]source}{}{}{}}
\makeatother
\crefname{originallemma}{original lemma}{original lemmas}
\Crefname{originallemma}{Original lemma}{Original lemmas}
\makeatletter
\newlabel{sites-cohomology-lemma-perfect}{{08G8}{source}{}{sourcepointer-7}{}}
\newlabel{sites-cohomology-lemma-perfect@cref}{{[originallemma][0][]08G8}{[1][0][]source}{}{}{}}
\makeatother
\crefname{originallemma}{original lemma}{original lemmas}
\Crefname{originallemma}{Original lemma}{Original lemmas}
\makeatletter
\newlabel{theorem-approximation}{{08HP}{source}{}{sourcepointer-8}{}}
\newlabel{theorem-approximation@cref}{{[originallemma][0][]08HP}{[1][0][]source}{}{}{}}
\makeatother
\title{Stacks, Derived Categories of Spaces: pseudo-coherent-hocolim}
\author{The Stacks Project Authors}
\date{stacks-project:spaces-perfect.tex:local-snapshot-20261012; GNU FDL 1.2 or later}
\begin{document}
\maketitle
\noindent\textit{Editorial note.} Statement span is lines 4536-4547 of sources/stacks-spaces-perfect/pinned-original-source.tex; proof span is lines 4549-4597, the complete author proof (47 body lines).\par\medskip
\noindent\textit{Editorial note.} Navigation stubs: 8 cited-but-not-included labels are supplied as non-load-bearing pointers with their pinned tags where the snapshot tag index has them.\par\medskip
\noindent\textit{Editorial note.} extra\_packages is empty on purpose: the retained preamble already loads xy with options and a second declaration clashes in the assembler wrapper.\par\medskip
\noindent\textit{Editorial note.} Text is retained verbatim, including the author's own typos and placeholders.\par\medskip
\noindent\textit{Editorial note.} Queue index 28 of queue\_spaceperfect.json (pb=1905, stmt 4536-4547, proof 4549-4597).\par\medskip
\noindent\textit{Editorial note.} stacks-spaces-perfect wave sp01 (2026-10-14): fresh adjudication accept per lane-s3/sp01\_judgments; pool per Chief 12:4x unfreeze (reserve-pool workflow, six-check clean); spaces-perfect.tex 126 with-proof rows, 70 b107-occupied labels avoided, 56 free rows judged 32 accept / 24 reject.\par\medskip
\noindent\textit{Editorial note.} Cross-chapter same-label warning (NOT an occupancy conflict, Chief ruling 2026-10-13 evening): label 'lemma-pseudo-coherent-hocolim' also exists in stacks-perfect — a different proposition in a different chapter file.\par\medskip
\noindent\textit{Allowed standard prerequisites (editorial).}\begin{itemize}
\item Explicit prerequisite boundary. The cited inputs are prerequisites; the argument in the retained span is not replaced by any cited result.
\end{itemize}\medskip
\ifdefined\hypertarget\hypertarget{sourcepointer-1}{}\fi
\noindent\textit{Source reference.} Original lemma 08HJ: sources/stacks-spaces-perfect/pinned-original-source.tex, source lines 3595--3605. Cited by the retained statement or proof but not part of the collected spans; supplied only as a navigation pointer. The argument does not depend on its content.\par\medskip
\ifdefined\hypertarget\hypertarget{sourcepointer-2}{}\fi
\noindent\textit{Source reference.} Original lemma 08HI: sources/stacks-spaces-perfect/pinned-original-source.tex, source lines 3575--3582. Cited by the retained statement or proof but not part of the collected spans; supplied only as a navigation pointer. The argument does not depend on its content.\par\medskip
\ifdefined\hypertarget\hypertarget{sourcepointer-3}{}\fi
\noindent\textit{Source reference.} Original lemma 0CRK: sources/stacks-spaces-perfect/pinned-original-source.tex, source lines 4597--4597. Cited by the retained statement or proof but not part of the collected spans; supplied only as a navigation pointer. The argument does not depend on its content.\par\medskip
\ifdefined\hypertarget\hypertarget{sourcepointer-4}{}\fi
\noindent\textit{Source reference.} Original lemma 0149: sources/stacks-spaces-perfect/pinned-original-source.tex, source lines 4597--4597. Cited by the retained statement or proof but not part of the collected spans; supplied only as a navigation pointer. The argument does not depend on its content.\par\medskip
\ifdefined\hypertarget\hypertarget{sourcepointer-5}{}\fi
\noindent\textit{Source reference.} Original lemma 09MB: sources/stacks-spaces-perfect/pinned-original-source.tex, source lines 4328--4338. Cited by the retained statement or proof but not part of the collected spans; supplied only as a navigation pointer. The argument does not depend on its content.\par\medskip
\ifdefined\hypertarget\hypertarget{sourcepointer-6}{}\fi
\noindent\textit{Source reference.} Original lemma 08FT: sources/stacks-spaces-perfect/pinned-original-source.tex, source lines 4597--4597. Cited by the retained statement or proof but not part of the collected spans; supplied only as a navigation pointer. The argument does not depend on its content.\par\medskip
\ifdefined\hypertarget\hypertarget{sourcepointer-7}{}\fi
\noindent\textit{Source reference.} Original lemma 08G8: sources/stacks-spaces-perfect/pinned-original-source.tex, source lines 4597--4597. Cited by the retained statement or proof but not part of the collected spans; supplied only as a navigation pointer. The argument does not depend on its content.\par\medskip
\ifdefined\hypertarget\hypertarget{sourcepointer-8}{}\fi
\noindent\textit{Source reference.} Original lemma 08HP: sources/stacks-spaces-perfect/pinned-original-source.tex, source lines 3786--3790. Cited by the retained statement or proof but not part of the collected spans; supplied only as a navigation pointer. The argument does not depend on its content.\par\medskip
\begin{lemma}
\label{lemma-pseudo-coherent-hocolim}
Let $S$ be a scheme.
Let $X$ be a quasi-compact and quasi-separated algebraic space over $S$.
Let $K \in D(\mathcal{O}_X)$. The following are equivalent
\begin{enumerate}
\item $K$ is pseudo-coherent, and
\item $K = \text{hocolim} K_n$ where
$K_n$ is perfect and $\tau_{\geq -n}K_n \to \tau_{\geq -n}K$
is an isomorphism for all $n$.
\end{enumerate}
\end{lemma}

\begin{proof}
The implication (2) $\Rightarrow$ (1) is true on any ringed site.
Namely, assume (2) holds. Recall that a perfect object of the derived
category is pseudo-coherent, see
Cohomology on Sites, Lemma \ref{sites-cohomology-lemma-perfect}.
Then it follows from the definitions that
$\tau_{\geq -n}K_n$ is $(-n + 1)$-pseudo-coherent
and hence $\tau_{\geq -n}K$ is $(-n + 1)$-pseudo-coherent,
hence $K$ is $(-n + 1)$-pseudo-coherent. This is true for
all $n$, hence $K$ is pseudo-coherent, see
Cohomology on Sites, Definition
\ref{sites-cohomology-definition-pseudo-coherent}.

\medskip\noindent
Assume (1). We start by choosing an approximation
$K_1 \to K$ of $(X, K, -2)$ by a perfect complex $K_1$, see
Definitions \ref{definition-approximation-holds} and
\ref{definition-approximation} and
Theorem \ref{theorem-approximation}.
Suppose by induction we have
$$
K_1 \to K_2 \to \ldots \to K_n \to K
$$
with $K_i$ perfect such that
such that $\tau_{\geq -i}K_i \to \tau_{\geq -i}K$ is an isomorphism
for all $1 \leq i \leq n$. Then we pick $a \leq b$ as in
Lemma \ref{lemma-ext-from-perfect-into-bounded-QCoh}
for the perfect object $K_n$. Choose an approximation
$K_{n + 1} \to K$ of $(X, K, \min(a - 1, -n - 1))$.
Choose a distinguished triangle
$$
K_{n + 1} \to K \to C \to K_{n + 1}[1]
$$
Then we see that $C \in D_\QCoh(\mathcal{O}_X)$ has
$H^i(C) = 0$ for $i \geq a$. Thus by our choice of $a, b$
we see that $\Hom_{D(\mathcal{O}_X)}(K_n, C) = 0$.
Hence the composition $K_n \to K \to C$ is zero. Hence by
Derived Categories, Lemma \ref{derived-lemma-representable-homological}
we can factor $K_n \to K$ through $K_{n + 1}$
proving the induction step.

\medskip\noindent
We still have to prove that $K = \text{hocolim} K_n$.
This follows by an application of
Derived Categories, Lemma \ref{derived-lemma-cohomology-of-hocolim}
to the functors
$H^i( - ) : D(\mathcal{O}_X) \to \textit{Mod}(\mathcal{O}_X)$
and our choice of $K_n$.
\end{proof}

\end{document}
